\chapter{System of relations}

	\noindent
	Let's put all the previous concepts in one table of data, to clarify what we have until now. What you are going to see is a mathematical miracle. TODAY you will say it is a simple table, in 20 years you will call it "beuatifull and simple, extremely elegant":
	\\\\
	
	\section{The table}
	
	\noindent
	\begin{table}[h!]
		\begin{center}
			%\begin{turn}{270}
			\begin{tabular}{|l|l|| l|l|l|l|l|l|l|l| l |l|l|l| l|}
				\hline
				\multicolumn{15}{|c|} {\cellcolor{lightgray} { } } \\
				\multicolumn{15}{|c|} {\cellcolor{lightgray} {\LARGE Totallity of ($SNEIs$ X $SNEIs$) } } \\
				\multicolumn{15}{|l|} {\cellcolor{lightgray} {$\:\:\:\:\:\:\:\:\:\:\:LCF_{2p}$ } } \\
				\hline
				\hline
				\cellcolor{lightgray} $r_{\theta_{k}}$ &
				\cellcolor{lightgray} $\theta_{k}$     &
				$F_{\infty}$ &
				$F_{0}$ &
				$F_{1}$ &
				$F_{2}$ &
				$F_{3}$ &
				$F_{4}$ &
				$F_{5}$ &
				$F_{6}$ &
				... &
				$F_{k-1}$
				& $F_{k}$
				&$F_{k+1}$
				&...  \\
				\hline
				\hline
				$r_{\theta_{1}}$ &
				$\theta_{1}$     &
				\cellcolor{applegreen!50}   &
				\cellcolor{applegreen!50}   &
				\cellcolor{bosunired!50}    &
				\cellcolor{bosunired!50}    &
				\cellcolor{bosunired!50}    &
				\cellcolor{bosunired!50}    &
				\cellcolor{bosunired!50}    &
				\cellcolor{bosunired!50}    &
				\cellcolor{bosunired!50}    &
				\cellcolor{bosunired!50}    &
				\cellcolor{bosunired!50}    &
				\cellcolor{bosunired!50}    &
				\cellcolor{bosunired!50}    \\
				\hline
				\hline
				$r_{\theta_{2}}$ &
				$\theta_{2}$     &
				\cellcolor{applegreen!50}   &
				\cellcolor{applegreen!50}   &
				\cellcolor{applegreen!50}   &
				\cellcolor{bosunired!50}    &
				\cellcolor{bosunired!50}    &
				\cellcolor{bosunired!50}    &
				\cellcolor{bosunired!50}    &
				\cellcolor{bosunired!50}    &
				\cellcolor{bosunired!50}    &
				\cellcolor{bosunired!50}    &
				\cellcolor{bosunired!50}    &
				\cellcolor{bosunired!50}    &
				\cellcolor{bosunired!50}    \\
				\hline
				\hline
				$r_{\theta_{3}}$ &
				$\theta_{3}$     &
				\cellcolor{applegreen!50}   &
				\cellcolor{applegreen!50}   &
				\cellcolor{applegreen!50}   &
				\cellcolor{applegreen!50}   &
				\cellcolor{bosunired!50}    &
				\cellcolor{bosunired!50}    &
				\cellcolor{bosunired!50}    &
				\cellcolor{bosunired!50}    &
				\cellcolor{bosunired!50}    &
				\cellcolor{bosunired!50}    &
				\cellcolor{bosunired!50}    &
				\cellcolor{bosunired!50}    &
				\cellcolor{bosunired!50}    \\
				\hline
				\hline
				$r_{\theta_{4}}$ &
				$\theta_{4}$     &
		        \cellcolor{applegreen!50}   &
				\cellcolor{applegreen!50}   &
				\cellcolor{applegreen!50}   &
				\cellcolor{applegreen!50}   &
				\cellcolor{applegreen!50}   &
				\cellcolor{bosunired!50}    &
				\cellcolor{bosunired!50}    &
				\cellcolor{bosunired!50}    &
				\cellcolor{bosunired!50}    &
				\cellcolor{bosunired!50}    &
				\cellcolor{bosunired!50}    &
				\cellcolor{bosunired!50}    &
				\cellcolor{bosunired!50}    \\
				\hline
				\hline
				$r_{\theta_{5}}$ &
				$\theta_{5}$     &
				\cellcolor{applegreen!50}   &
				\cellcolor{applegreen!50}   &
				\cellcolor{applegreen!50}   &
				\cellcolor{applegreen!50}   &
				\cellcolor{applegreen!50}   &
				\cellcolor{applegreen!50}   &
				\cellcolor{bosunired!50}    &
				\cellcolor{bosunired!50}    &
				\cellcolor{bosunired!50}    &
				\cellcolor{bosunired!50}    &
				\cellcolor{bosunired!50}    &
				\cellcolor{bosunired!50}    &
				\cellcolor{bosunired!50}    \\
				\hline
				\hline
				$r_{\theta_{6}}$ &
				$\theta_{6}$     &
				\cellcolor{applegreen!50}   &
				\cellcolor{applegreen!50}   &
				\cellcolor{applegreen!50}   &
				\cellcolor{applegreen!50}   &
				\cellcolor{applegreen!50}   &
				\cellcolor{applegreen!50}   &
				\cellcolor{applegreen!50}   &
				\cellcolor{bosunired!50}    &
				\cellcolor{bosunired!50}    &
				\cellcolor{bosunired!50}    &
				\cellcolor{bosunired!50}    &
				\cellcolor{bosunired!50}    &
				\cellcolor{bosunired!50}    \\
				\hline
				\hline
				$r_{\theta_{7}}$ &
				$\theta_{7}$     &
				\cellcolor{applegreen!50}   &
				\cellcolor{applegreen!50}   &
				\cellcolor{applegreen!50}   &
				\cellcolor{applegreen!50}   &
				\cellcolor{applegreen!50}   &
				\cellcolor{applegreen!50}   &
				\cellcolor{applegreen!50}   &
				\cellcolor{applegreen!50}   &
				\cellcolor{bosunired!50}    &
				\cellcolor{bosunired!50}    &
				\cellcolor{bosunired!50}    &
				\cellcolor{bosunired!50}    &
				\cellcolor{bosunired!50}    \\
				\hline
				\hline
				... &
				... &
				\cellcolor{applegreen!50}   &
				\cellcolor{applegreen!50}   &
				\cellcolor{applegreen!50}   &
				\cellcolor{applegreen!50}   &
				\cellcolor{applegreen!50}   &
				\cellcolor{applegreen!50}   &
				\cellcolor{applegreen!50}   &
				\cellcolor{applegreen!50}   &
				\cellcolor{applegreen!50}   &
				\cellcolor{bosunired!50}    &
				\cellcolor{bosunired!50}    &
				\cellcolor{bosunired!50}    &
				\cellcolor{bosunired!50}    \\
				\hline
				\hline
				$r_{\theta_{k}}$ &
				$\theta_{k}$     &
				\cellcolor{applegreen!50}   &
				\cellcolor{applegreen!50}   &
				\cellcolor{applegreen!50}   &
				\cellcolor{applegreen!50}   &
				\cellcolor{applegreen!50}   &
				\cellcolor{applegreen!50}   &
				\cellcolor{applegreen!50}   &
				\cellcolor{applegreen!50}   &
				\cellcolor{applegreen!50}   &
				\cellcolor{applegreen!50}   &
				\cellcolor{bosunired!50}    &
				\cellcolor{bosunired!50}    &
				\cellcolor{bosunired!50}    \\
				\hline
				\hline
				$r_{\theta_{k+1}}$ &
				$\theta_{k+1}$     &
				\cellcolor{applegreen!50}   &
				\cellcolor{applegreen!50}   &
				\cellcolor{applegreen!50}   &
				\cellcolor{applegreen!50}   &
				\cellcolor{applegreen!50}   &
				\cellcolor{applegreen!50}   &
				\cellcolor{applegreen!50}   &
				\cellcolor{applegreen!50}   &
				\cellcolor{applegreen!50}   &
				\cellcolor{applegreen!50}   &
				\cellcolor{applegreen!50}   &
				\cellcolor{bosunired!50}    &
				\cellcolor{bosunired!50}    \\
				\hline
				\hline
				$r_{\theta_{k+2}}$ &
				$\theta_{k+2}$     &
				\cellcolor{applegreen!50}   &
				\cellcolor{applegreen!50}   &
				\cellcolor{applegreen!50}   &
				\cellcolor{applegreen!50}   &
				\cellcolor{applegreen!50}   &
				\cellcolor{applegreen!50}   &
				\cellcolor{applegreen!50}   &
				\cellcolor{applegreen!50}   &
				\cellcolor{applegreen!50}   &
				\cellcolor{applegreen!50}   &
				\cellcolor{applegreen!50}   &
				\cellcolor{applegreen!50}   &
				\cellcolor{bosunired!50}    \\
				\hline
				\hline
				... &
				... &
				\cellcolor{applegreen!50}   &
				\cellcolor{applegreen!50}   &
				\cellcolor{applegreen!50}   &
				\cellcolor{applegreen!50}   &
				\cellcolor{applegreen!50}   &
				\cellcolor{applegreen!50}   &
				\cellcolor{applegreen!50}   &
				\cellcolor{applegreen!50}   &
				\cellcolor{applegreen!50}   &
				\cellcolor{applegreen!50}   &
				\cellcolor{applegreen!50}   &
				\cellcolor{applegreen!50}   &
				\cellcolor{applegreen!50}   \\
				\hline
				
			\end{tabular}
			\caption{System of relations $r_{\theta_{k}}$}
			\label{tab: system_relations_r_theta_k}
			%\end{turn}
		\end{center}
	\end{table}
	\\\\
	
	\newpage
	
	\section{Observations:}
	
	\subsection{Description of the table:}
	
	\noindent
	GREEN: Means that relation is able to offer COMPLETELY disjoint PACKS for those concrete D-Pairs inside those Families.\\\\
	RED: Means PACKs assigned have some elements in common. The PACKS are not "clean".\\\\
	K is a natural number.\\
	\\\\
	
	\subsection{Derivations of general rule:}
	
	\noindent
	We can observe clearly the general rule:\\
	\fbox{
		\textbf{\textit{$r_{\theta_{k}}$ is able to solve ALL D-Pairs inside the Families: $F_{k-1},\: F_{k-2}, \:, \:..\: , \: F_{0} ,\: F_{\infty}$.}}
	}
	\\\\
	
	\noindent
	It means too:\\\\
	1) WE have covered ALL ( SNEIs X SNEIs ): for each possible D-pair, it belongs to a unique Family of D-pairs, ALL with the same $\gamma$-value, ALL unsolved in a FINITE quantity of relations $r_{\theta_{k}}$, BUT SOLVED in an INFINITE quantity of relations $r_{\theta_{k}}$.\\\\
	2) From the spectrum of just D-pairs, WE ARE INMUNE TO ALL POSSIBLE DIAGONALIZATIONS. You can not mention a single D-pair, not solved in infinite relations. ALL relations PREVIOUSLY defined, until the last R-pair, BEFORE you have mentioned the D-pair.\\\\
	3) The table is like a function by infinite parts. Each universe $\theta_{k}$, following a concrete sub-relation, is able to "solve" a new subset of D-pairs. If we quit the increasing capacity of solving D-pairs... from the perspective of D-pairs, it is LIKE a function by parts, where each subset, solve one subset from the other set. IT IS WEIRD, because it is not directly the spectrum of single elements... but this is a clear CLUE $\aleph_{1}$ is not untouchable. And...\\\\
	4) A CLEAR PROOF that the spectrum of elements and D-pairs are not equivalents. Not always what happens in the spectrum of single elements, happens in the spectrum of D-pairs... D-pairs are TOTALLY SOLVABLE, using all relations, each one uses a disjoint subset of $\mathbb{N}$. This is a crack in definitions.\\\\
	\\\\
	
	\noindent
	These are DIRECT observations. Based in large, but simple constructions. I have not used complicated mathematical tools.
	\\\\
	
	\subsection{Comments}
	
	\noindent
	1) The trick to shoot arrows to the "paradise" Cantor created for us, was shooting more than one at the same time. Like an army, splitting itself in different lines of defense, each one with a different strategy. Exactly, we needed $\aleph_{0}$ different arrows, different parallel relations. Counting naturals in order, transforming them into members of LCF, and seeing if they belong to $LCF_{2p}$, or not, we use less than 5\% of natural numbers... until 300 millions. And the quantity tends to decrease, observing the graphic: "x" is natural number; "y" is the quantity of members inside $LCF_{2p}$, $LCF_{1}$ or Previous, divided by the quantity of ALL natural numbers, until that point in x. Usefull ones divided by total ones. The difference is inside $LCF_{2c}$, more than 95\% of $\mathbb{N}$ until 300 millions.
	\\\\
	
	\noindent
	2) $f(a) \:=\: f(b) \:\leftrightarrow\: a \:=\: b $: The core of the definition of injectivity. We have a mathematical FACT: we can solve all D-pairs, but we can not solve the comparison from the perspective of single elements ( Cantor's diagonalization ). IT HAPPENS. This construction has no fails, and is very simple. The problem with the definition is, that it lets me focus ONLY in D-pairs... and that the system of relations is WORST than a function by parts. Each universe INCREASES ITS HABILITY TO SOLVE D-PAIRS, they do not only solve one Family.
	
	\noindent
	If we put a perfect perpendicular light, comming from the bottom of the table, that were, in the infinity... THE ENTIRE LINE OF FAMILIES WOULD BE "SHADOWED" by the pattern of the semi-xmas-tree, of solved Families, in each relation $r_{\theta_{k}}$. Our simple subsets of $\mathbb{N}$, grows, from an undeterminated state of being able to solve only TWO Families of D-pairs, to solve nothing less than ALL Families of D-pairs. All Families means ALL (SNEIs X SNEIs). And solving ALL the cartessian product is an equivalent of solving all the set SNEIs, that has the same cardinal than $\mathcal{P}(\mathbb{N})$. This can be said just using the definition of injectivity.
	\\\\
	
	\noindent
	3)This is one of the greatest Egg of Columbus if mathematics history. (SNEIs X SNEIs) is a set more complex than only SNEIs. If you say people, you splitted a set with cardinal $\aleph_{1}$, into $\aleph_{0}$ disjoint subsets ( Families ), only quick people with talent see inmediately that some of those subsets MUST HAVE cardinal $\aleph_{1}$. If you say you can "solve" each one, using for each one, a different subset from a partition of $\mathbb{N}$... wer are saying "\textbf{\textit{we can solve a set with cardinal $\aleph_{1}$ only using a set with cardinal $\aleph_{0}$}}". And they overreact and refuse to hear, THIS large, but simple solution: we haven't used complicated tools... and the core depends only in positions of symbols, and well choosen combinatorics.

	\noindent
	The other group of experts, that only see step by step, in order, without any advice about what we are going to do... only are able to see, we solve Families entirely and  correctly. But they consider all this as boring... and without value.
	\\\\
	
	\noindent
	The solution did not cost them 20 years as it cost to me. And it only solves what others experts called "mathematically impossible": an Egg of Columbus. We solve an IMPOSSIBLE problem, with an insultingly simple solution. A "boring one".
	\\\\


	\noindent
	4) We have another weird effect: solving ALL D-pairs, MUST mean we should be able to solve all single elements, from the spectrum of single elements and functions, not relations. But we can't. At the same time, not being able to solve the spectrum of single elements ( the core of Cantor's proof) MUST mean we should NOT be able to solve all D-pairs with previoulsy defined relations: but we can.
	
	\noindent
	This is a new mathematical phenomenom ignored by Cantor. And he forgot to reflect these posibilities in his definitions. This is not a contradiction: this is a buildable phenomenom. It can contradict OTHER theorems, ideas or definitons, but itself, is not an absurd. And it is easy to be builded.
	\\\\
	
	
	\noindent
	5) Until here, we have one third of the entire presentation. The first phase of three ( without counting the final fourth one ). This is my first "mathematically impossible" intermediate point. And it has value by itself. It is an AMAZING, and INTERESTING mathematical phenomenom, by itself. WE covered entirely a set with cardinal $\aleph_{1}$. But we did it in a very weird way... "not so clear"... 
	
	\noindent
	Only with this, you should run to try to contact me and talk about ONLY THIS. WAIT, NOT TALK, directly to do conferences in your university. But I know you are in the mood: "trying to see where this crankery commited a mistake". And you forgot to stop a little, and just simply be amazed by the beauty of this phenomenom isolated. I need to make you confess. In case you are not convinced, don't worry. Or better said: worry. The worst is coming. The next two phases are designed, based on this table, to make you GUARANTEE, a practical equivalent of an injectivity $r:\mathcal{P}(\mathbb{N}) \longrightarrow \mathbb{N}$ exists !!
	\\\\
	
	\noindent
	You need to read the next TWO phases. Not only one and getting angry a lot. PRECISELY you are going to GUARANTEE that equivalent is buildable, when you were thinking you have found the mistake of one phase ( or branch ), ignoring you are guaranteeing the weak point of the other one, "while" you are trying to describe the mistake. And viceversa. YOU and not me, are going to guarantee it. YOU are the expert, not me.
	\\\\
	
	\noindent
	OBLIGATORY READ: unnecesary comment 006. The Shit-tuation. If you haven't read it until now, read it NOW, before continuing. If you have read it, read it again... you will compare it to the table... but it is not THAT exactly, what we are going to do in DI and TPI... but you will have an idea about HOW it could happen.
	\\\\
	
	\noindent
	END OF PHASE 001.
	